\section{実在する混合$A_2$と二段拡大の非消失}\label{sec:dp-hearts2-mixed-a2}
本節の$E_6$最大共鳴の整数表示とZariski稠密性は主論文\cite{Hearts}を入力とする．
数値的共鳴だけから任意の根sheafの存在を推測せず，指定した単純根の実在を確認する．

\begin{proposition}[実在する正負の単純根sheaf]
\label{prop:actual-mixed-a2}
\[
 (N_1,\ldots,N_6)=(1,3,2,5,7,11),\qquad p_i=N_iq
\]
とする．六単純根のEuler指数は
\[
 (m_0,m_1,m_2,m_3,m_4,m_5)=(6,2,-1,3,2,4).
\]
特に
\[
 \alpha=e_1-e_2,\qquad\beta=e_2-e_3,\qquad
 m(\alpha)=2,\quad m(\beta)=-1,\quad
 \alpha\cdot\beta=1.
\]
Rains\cite[\S4, (4.14)--(4.15)]{RainsFiltered}の
simple-root effectivity判定と
\cite[\S9,\S11]{RainsBirational}の既約sheafの存在を
このsmooth elliptic chartに適用すると，同じRains heartに
境界非交叉の既約root sheaf
$M_\alpha,M_\beta$が存在する．それらについて
\[
 \dim\Ext^1(M_\alpha,M_\beta)=1,\quad
 \dim V^0(M_\alpha)=2,\quad
 \dim V^1(M_\beta)=1.
\]
またこの性質を保ち$N_1<N_3<N_2$とする全根共鳴点の集合も
同じ六点配置chartでZariski稠密である．
\end{proposition}
\begin{proof}
単純根指数は$h|_E=3o$と点の差から計算される．
全ての単純根の制限は$\langle q\rangle$にあるので，
Rainsのeffectivity判定によりこれらは有効である．
$E$は滑らかで$\alpha\cdot E=0$なので，
有効な単純根はformal $(-2)$-curveとして
境界と交わらない既約純一次元sheafで実現される．
無限位数の$q$のため$M_\alpha,M_\beta$のEuler指数は
上記の値に一意に固定される．
前節の命題\ref{prop:root-intersection}および
定理\ref{thm:sign}からExtと$V^j$の次元が従う．
最後は主論文\cite{Hearts}の全次数の最大共鳴と整数順序の稠密性による．
\end{proof}

\begin{proposition}[二段混合拡大全体での非消失]
\label{prop:a2-directsum}
命題\ref{prop:actual-mixed-a2}の二sheafを固定する．
$a\ge1,b\ge0$とし，任意の拡大
\[
0\longrightarrow M_\beta^{\oplus b}\longrightarrow M
 \longrightarrow M_\alpha^{\oplus a}\longrightarrow0
\]
を取る．$\Ext^1(M_\alpha,M_\beta)\simeq k$の生成元が
定める$A=R/(g-1)$加群写像を
$d:V_\alpha\to V_\beta$とする．
拡大類に対応する行列$T\in\operatorname{Mat}_{b\times a}(k)$の
connecting morphismは$d\otimes T$であり，
$r=\operatorname{rank}(d)\operatorname{rank}(T)$と置けば
\[
\boxed{
\dim V^0(M)=2a-r,\quad
\dim V^1(M)=b-r,\quad
r\le\min(a,b).}
\]
特に$a\ge1$なら$\dim V^0(M)\ge a>0$なので
自然な反標準導来切断はこの二段拡大を消さない．
$b=2a$では$\chi(e,M)=0$だが，依然として
$\Phi(M)\ne0$である．
\end{proposition}
\begin{proof}
直和とYoneda積の自然性により
\[
\Ext^1(M_\alpha^{\oplus a},M_\beta^{\oplus b})
 \cong\operatorname{Mat}_{b\times a}(k),
\]
対応するconnecting morphismは$k$-線型空間上
Kronecker積$d\otimes T$である．そのrankは積となり，
$d$は二次元から一次元への射だから$\operatorname{rank}(d)\le1$．
前節の六項完全列
\eqref{eq:offboundary-six}に代入して次元公式を得る．
\end{proof}

\begin{corollary}[全$q$導来同値に対する明示的な検査条件]
\label{cor:a2-ff}
同じ例で自然な導来切断$\Phi$が
$M_\alpha,M_\beta$の間の$\Ext^1$で完全忠実なら
\[
 \dim\Hom_A(V_\alpha,V_\beta)=1
\]
である．従って二次元表現$V_\alpha$は
一次元表現$V_\beta$を商に持ち，単純$A$加群ではない．
そのとき非分裂拡大のconnecting morphismはrank$1$で，
$(\dim V^0(M),\dim V^1(M))=(1,0)$となる．
\end{corollary}
\begin{proof}
前節の系\ref{cor:affine-ext-test}による．
$\dim V_\beta=1$だから非零$A$準同型は全射である．
$\Ext^1$の生成元は完全忠実性の下で非零の
$\Hom_A$へ写るから，connecting morphismのrankは$1$．
\end{proof}

\begin{remark}[有理数上の検算例]
\label{rem:a2-rational-witness}
\[
 E:\ y^2z=x^3-xz^2+z^3,\quad o=(0:1:0),\quad
 q=(0:1:1)
\]
とする．楕円曲線の二倍公式で
$2q=(\frac14:-\frac78:1)$である．
Nagell--Lutzの定理により$q$は無限位数である：
有限位数なら$2q$も有理torsionで整数座標を持つはずだが
$\frac14$は整数でない．
$N=(1,3,2,5,7,11)$の六点について
20通りの三点行列式と六点のconic評価行列式が
全て非零であり，72根の共鳴指数も全て非零であることを
\texttt{computations/verify\_resonance\_maximal\_a2.py}
により\emph{有理数の厳密演算}で検算した．
このコードは$A$の表現行列や$d$のrankを
計算するものではなく，Rainsのsheaf存在を
数値実験で置き換えてはいない．
\end{remark}
